% TOP_PROBLEM: {"id": 3207, "title": "Pentagon problem", "category_id": 3, "category": {"id": 3, "name": "graph_theory", "display_name": "Graph Theory", "description": "Problems involving graphs, networks, and their properties.", "slug": "graph-theory", "order_index": 3, "created_at": "2026-07-31T15:26:25.671Z"}, "status": "open", "problem_number": "OPG-167", "set_id": 12}
% FIRST_POSTED: 2026-10-01
% ENTRY_KIND: statement_only
% UnsolvedMath ID 3207; source catalogue code OPG-167
% !TeX program = lualatex
% Definition and sources only; this file is not an LLM solution attempt.
\documentclass[11pt,a4paper]{article}
\usepackage[margin=25mm]{geometry}
\usepackage{fontspec}
\setmainfont{lmroman10-regular.otf}[BoldFont=lmroman10-bold.otf,
 ItalicFont=lmroman10-italic.otf,BoldItalicFont=lmroman10-bolditalic.otf]
\usepackage{amsmath,amssymb,mathtools}
\usepackage{unicode-math}
\setmathfont{latinmodern-math.otf}
\AtBeginDocument{\let\setminus\smallsetminus}
\usepackage{xurl,hyperref}
\hypersetup{colorlinks=true,urlcolor=blue}
\setcounter{secnumdepth}{0}
\setlength{\parindent}{0pt}
\setlength{\parskip}{5pt}
\setlength{\emergencystretch}{3em}
\begin{document}
\section*{Pentagon problem}\label{3207}
{\small UnsolvedMath ID 3207 · OPG-167 · Graph Theory · OpenGarden}
\subsection{Definitions and mathematical statement}
Let $G$ be a finite simple cubic graph, so each
vertex has exactly three neighbours. Its girth
is the length of its shortest simple cycle.
Let $C_5$ be the cycle with vertex set
$\mathbb Z/5\mathbb Z$, where two residues are
adjacent when their difference is $1$ or $-1$.
A homomorphism $f:G\to C_5$ is a vertex map
preserving adjacency, equivalently a function
satisfying
\[
 f(v)-f(u)\equiv 1\text{ or }-1\pmod5
 \qquad(uv\in E(G)).
\]
The map need not be injective or surjective.

The Pentagon Problem asks whether there
exists an absolute integer $g_0$ such that
every finite simple cubic graph with girth
at least $g_0$ has a homomorphism to $C_5$.
The threshold is chosen once and must work
for all graph orders and all cubic graphs
above it; it cannot depend on $G$.

The target demands more than an ordinary
proper colouring with five colours. Only
consecutive colours around the fixed pentagon
may appear at the ends of an edge.
Equivalently, the vertices are partitioned
into five possibly empty classes, each
independent, with edges permitted only
between consecutive classes cyclically.
The hypothesis forbids all short cycles,
not only short odd cycles. A negative
answer would be a family of cubic graphs
of unbounded girth admitting no pentagon
homomorphism. Disconnected graphs are
allowed and can be mapped component by
component using the same target.
\subsection{Short English statement}
Does one sufficiently large girth threshold guarantee that every cubic graph admits an adjacency-preserving map to a pentagon?
\subsection{Sources}
\begin{itemize}
\item[S1] ULAM.AI and the UnsolvedMath Contributors, supplied problems.json, record 3207 (OPG-167), OpenGarden collection. Catalogue identity and source transcription; CC BY 4.0.
\url{https://huggingface.co/datasets/ulamai/UnsolvedMath}.
\item[S2] Open Problem Garden, Pentagon problem. Original problem and discussion; the mathematical formulation below is expanded from the supplied source record.
\url{http://www.openproblemgarden.org/op/pentagon_problem}.
\end{itemize}
\end{document}
